\chapter{Integrating Inferred Specifications into the AI-Native Lifecycle}\label{ch:integration}

The two preceding chapters established that specifications can be inferred cheaply
and put to use. \Cref{ch:inferrer} showed that heuristic inference recovers sound,
verifiable contracts from Java source and that supplying them to a language model
measurably improves the tests it generates; \cref{ch:compositional} showed that a
second, weakest-precondition pass strictly extends what single-pass inference can
recover and, in doing so, turns behavioural version compatibility into a static
analysis. Each treated an inferred contract as an end in itself---an oracle for
test generation, a property to refine. This chapter draws the threads into a
single lifecycle argument. It asks what a software development process would look
like if it were reorganised around a machine-checkable contract as its durable
source of truth, and answers with a verification-centric \emph{reference
architecture} in which that contract is the pivot between ambiguous human intent
and precise machine verification. The instruments of the earlier chapters reappear
here not as fresh contributions but as the trusted components from which the
architecture is assembled. The conceptual background---program synthesis, the
test-oracle problem, and behavioural version compatibility---is surveyed in
\cref{ch:background} and not repeated.

The motivation is an economic inversion already argued in \cref{ch:introduction}.
When an agent can regenerate a conforming implementation from a description in
seconds, the implementation ceases to be the treasured asset and becomes closer to
a disposable build output; what becomes scarce and durable is the specification of
intent against which any number of cheap regenerations are measured. The emerging
agentic, specification-driven lifecycles---GitHub's Spec Kit~\citep{githubspeckit},
Amazon's AI-Driven Development Life Cycle~\citep{awsaidlc}, and the Kiro
integrated development environment~\citep{awskiro}---have noticed the inversion and
place a ``specification'' at their centre, but the specification they organise
around is natural language, which is not machine-checkable. Examined as a control
system, these lifecycles exhibit three compounding defects: the spec-to-code step
is an unverifiable compilation from an informal source; the loop is open, closed
only by human review or a sampled test suite; and where automated acceptance
exists it is circular, because the same model that writes the code also writes the
oracle it is graded against~\citep{liu2023evalplus,liu2026oracle}. These are one
missing component viewed from three angles: a \emph{verification spine}. This
chapter supplies it.

\section{Enabling machinery: carriage and continuous integration}\label{sec:int-enablers}

Before the architecture can be described, two prior results must be recalled as
enablers, because a lifecycle organised around specifications presupposes that
specifications can travel with the artefacts developers exchange and can be
produced routinely without disrupting the team. Both were established as
contributions in their own right; here they are compressed to the role they play
in the larger pipeline.

\subsection{Carriage: specifications that survive compilation and publication}\label{sec:int-carriage}

A specification that cannot travel with the artefact it describes lives in a
separate document that drifts out of date and is unavailable to the very
consumers---verifiers, IDEs, test generators, and models reading an API---that
would benefit from it. Java libraries ship as JAR files and REST endpoints expose
only an HTTP signature, so source is the exception rather than the norm. The
carriage mechanism resolves this by writing inferred JML~\citep{leavens2006jml}
into runtime-retained class-file annotations via an ASM-based embedder, with a
parallel OpenAPI~3.x extension family (\texttt{x-jml-*}) for service boundaries and
a JML stub-file sidecar that OpenJML~\citep{cok2014openjml} reads natively for
verification. The design descends from the design-by-contract observation that a
contract is a property of the supplier as seen by a client who most often holds
only the binary~\citep{meyer1992dbc}. Empirically it round-trips losslessly:
across six mature libraries totalling 39{,}268 methods, and across 22{,}881 real
inferred specifications recovered from those libraries' source, every specification
the harness embeds is recovered byte-for-byte, at a bytecode overhead of roughly
4.6--10.2\% after format optimisation and at throughputs that comfortably support
per-pull-request continuous integration. Embedded classes load correctly in JVMs
that have never seen the annotation type, so a JAR can be enriched in place with no
risk to consumers and no toolchain change. The single fact the lifecycle needs is
that a contract produced by inference can be carried, at full fidelity, on the
artefact a developer actually deploys.

\subsection{Continuous integration: inference as a routine, fail-open step}\label{sec:int-ci}

For inference to change how software is built it must run where development
happens---in the CI pipeline, on every change, without disrupting the merge train.
Formal-method tooling is conventionally held to be too slow and too brittle for
this, and three deliberate design choices overturn that calculus. First,
\emph{diff-only analysis with content-hash caching}: whole-codebase re-inference
runs once per merge to the trunk, while per-pull-request work is restricted to
changed methods, unchanged methods being served from cache. The cache key is a
SHA-256 hash of the method source together with the canonical form of its callees'
specifications and the inferrer version, so a callee's specification change
invalidates precisely its callers and nothing else---the property that makes the
cache friendly to the compositional analysis of \cref{ch:compositional}. Second,
\emph{fail-open semantics}: an inference crash, an OpenJML timeout, or a malformed
clause is logged to the pull-request comment but never blocks the merge, so the
pipeline is safe to deploy where formal methods have previously been resisted;
strict gating is available only as an opt-in. Third, \emph{persistent in-place
reporting}: a reviewer sees one collapsible summary per pull request rather than a
growing thread of comments. A working reference implementation---a command-line
driver with \texttt{infer}, \texttt{verify}, \texttt{embed}, and \texttt{summary}
subcommands, plus templates for GitHub Actions, GitLab CI, and Jenkins---is
dogfooded against the inferrer's own repository, where a clean-cache run completes
in roughly ninety seconds and a deliberately introduced syntax error is surfaced
without blocking the merge. The controlled study of pipeline overhead across
external repositories is designed but not yet run, a boundary \cref{ch:discussion}
revisits. The enabling fact is that inference can be a routine CI citizen: cheap,
incremental, and unobtrusive enough to run as a matter of course.

\section{A verification-centric reference pipeline}\label{sec:int-pipeline}

With carriage and CI in place, the lifecycle can be reorganised around a
machine-checkable contract. The pipeline is offered as a reference
architecture---a set of design principles, a layered model, and an end-to-end data
flow---validated by instantiation rather than asserted as a proven methodology.
\Cref{fig:int-pipeline} shows its shape.

\begin{figure}[t]
  \centering
  \begin{tikzpicture}[
    font=\small,
    box/.style={draw, rounded corners, align=center, text width=7.4cm,
                minimum height=0.85cm, fill=white, inner sep=3pt},
    pivot/.style={draw, very thick, align=center, text width=7.9cm,
                  minimum height=0.85cm, fill=blue!10, inner sep=3pt},
    lbl/.style={font=\itshape\footnotesize, align=left},
    arr/.style={-{Latex[length=2.2mm]}, semithick},
  ]
    \node[box] (intent) {\textbf{Intent} \footnotesize --- human goals, tickets, docs, existing code};
    \node[box, below=0.7cm of intent] (elab)
      {\textbf{Elaboration} \footnotesize(elaborator role): ratify \emph{behaviour}
       by example; adversarial + vacuity checks};
    \node[pivot, below=0.7cm of elab] (contract)
      {\textbf{Formal contract} --- the \textsc{pivot} \footnotesize (a set of
       $\langle$provenance, assurance$\rangle$-tagged obligations)};
    \node[box, below=0.7cm of contract] (synth)
      {\textbf{Synthesis} \footnotesize(implementer role): contract $\rightarrow$
       code, a disposable build output};
    \node[box, below=0.7cm of synth] (verif)
      {\textbf{Verification} \footnotesize: code $\models$ contract? OpenJML ESC
       $+$ Z3, over \emph{all} inputs};
    \draw[arr] (intent) -- (elab);
    \draw[arr] (elab) -- (contract);
    \draw[arr] (contract) -- node[right=2pt,lbl]{contract as \emph{input}} (synth);
    \draw[arr] (synth) -- (verif);
    \draw[arr] (verif.east) -- ++(0.5,0) |- ([yshift=2pt]synth.east)
      node[pos=0.25,right=1pt,lbl,text width=2.5cm]{fail: counterexample (concrete witness)};
    \begin{scope}[on background layer]
      \node[fill=orange!10, rounded corners, fit=(intent)(elab), inner sep=6pt,
            label={[font=\itshape\footnotesize]above:ambiguous, human-owned world}] {};
      \node[fill=green!10, rounded corners, fit=(synth)(verif), inner sep=6pt,
            label={[font=\itshape\footnotesize]below:precise, machine-checkable world}] {};
    \end{scope}
  \end{tikzpicture}
  \caption{The reference pipeline. The formal contract is the pivot between the
  ambiguous, human-owned world and the precise, machine-checkable world.
  Elaboration converts intent into a contract under human ratification of
  behaviour; synthesis and verification form a counterexample-guided loop below the
  pivot; and the verifier never grades a contract it authored.}
  \label{fig:int-pipeline}
\end{figure}

\subsection{The contract as pivot, and its dual role}\label{sec:int-pivot}

A machine-checkable formal contract sits at the centre of the architecture as the
\emph{pivot} between the informal activity of deciding what is wanted and the
mechanically checkable activity of producing and certifying code. It is the single
point at which informality is converted into formality and the only artefact that
crosses that boundary. Concretely the pivot is expressed in
JML~\citep{leavens2006jml}, whose pre/postcondition discipline descends from Hoare
logic~\citep{hoare1969axiomatic}, Dijkstra's weakest-precondition
calculus~\citep{dijkstra1976discipline}, and Meyer's design by
contract~\citep{meyer1992dbc}, and discharged by OpenJML extended static
checking~\citep{cok2014openjml} over Z3~\citep{demoura2008z3}.

The decisive property is that the contract plays two roles at once. It is the
\emph{input}---the precise intent the synthesis agent builds against---and it is
the \emph{oracle}---the artefact against which that agent's output is formally
verified. This duality makes the pipeline self-consistent: there is exactly one
statement of intent, and code is judged against the very artefact that produced it,
which is precisely what the natural-language lifecycles cannot achieve, because
their input has no formal semantics against which output could be checked. The
historical obstacle to the dual role was the cost of authoring formal contracts;
that obstacle is removed by the inference of \cref{ch:inferrer} and
\cref{ch:compositional}, which supplies a draft automatically, so the human's task
collapses from \emph{authoring} formal text to \emph{ratifying} behaviour.

\subsection{A closed synthesis--verification loop}\label{sec:int-loop}

Below the pivot the ambiguous human world has been left behind: the contract is a
precise object, and the agent's task is no longer to interpret intent but to
satisfy an obligation. Here the open generate-then-eyeball cycle is replaced by a
closed loop shaped like counterexample-guided inductive synthesis
(CEGIS)~\citep{solarlezama2006sketching,solarlezama2008sketching}. A synthesis
agent proposes a body; the trusted verifier attempts to prove it satisfies the
contract over all inputs; and any failure is returned to the agent not as a verdict
but as a concrete counterexample---a fully instantiated pre-state in which the body
cannot discharge the obligation---that drives the next attempt. The value lies
entirely in the quality of that feedback. Self-improving code agents close their
loop with a sampled test suite, which certifies only the finitely many inputs it
happened to exercise and whose weakness as a correctness oracle is well
documented~\citep{liu2023evalplus,chen2021codex}; a failing test reports that one
input misbehaved but gives no assurance that fixing it will not break another. A
successful ESC proof, by contrast, certifies the obligation over the entire input
domain, and a failed one hands the synthesiser a witnessed cause in the same logic
the verifier reasons in. Because any method with a loop needs an inductive
invariant that ESC cannot guess, the loop separates two hard problems: the agent
contributes the algorithmic content, and the inferrer of \cref{ch:inferrer}
supplies the discharging loop invariant---in the direction-aware, inductively
preserved form ESC requires---for exactly the loop families it handles. When the
solver neither proves nor refutes within its budget---Z3 is incomplete on
multiplicative arithmetic, array theory, and richly quantified preconditions---the
obligation is not dropped but \emph{degraded} down a ladder from full proof to
bounded checking to a generated test, with the level achieved recorded on the
obligation rather than silently lowered. This graduation is what makes a
verification-centric lifecycle realistic rather than aspirational.

\subsection{Independent authority: a model must not grade its own work}\label{sec:int-authority}

Because the contract is the oracle, verification is meaningful only if the contract
is not authored by the same agent that writes the code. If one model both proposes
the contract and satisfies it, the verifier still soundly proves that the code
meets the contract, but that proof certifies nothing about \emph{intent}: a shared
misreading is encoded identically in contract and code, which then agree without
exposing it. This is the self-grading circularity that undermines test-suite
acceptance and, more acutely, model-driven oracle
generation~\citep{liu2026oracle}. The \emph{independent-authority constraint}
therefore partitions the lifecycle into an \emph{elaborator} role---maximally
faithful, never permitted to see the implementation, and the owner of the
ratified examples---and an \emph{implementer} role, whose only task is to satisfy
the contract handed to it. The verifier, reused unchanged, is authored by neither.
The constraint's scope must be stated honestly: it removes self-grading, but it
does not by itself decorrelate a shared misreading of intent across two roles of
the same underlying model. That residual is broken only by human ratification,
whose authority lies outside the model distribution. Ratification is made feasible
by a discipline of \emph{ratifying behaviour, not logic}: rather than confirm a
page of quantified JML, the human is asked concrete questions in domain
terms---``given an empty list, this contract throws rather than returns zero; is
that what you want?''---and each confirmed answer becomes a machine-checkable
pinned example the contract must henceforth satisfy. Every obligation carries a
provenance tag recording its origin and an assurance tag recording the strongest
verification level achieved, so that adherence to the constraint is auditable and
the word ``verified'' stays honest: a reader can see which obligations are
human-ratified, which are merely machine-inferred, and which were proved as opposed
to merely tested.

\subsection{Brownfield contract recovery by triangulation}\label{sec:int-brownfield}

Most software is not new, so in the common case no contract exists and one must be
\emph{recovered}. The difficulty is that the existing code is at once the best
evidence of what the system does and the prime suspect, since any bug it contains
is encoded in exactly that code; a contract inferred from code alone therefore
describes the implementation faithfully, defects included. The recovery phase
consequently triangulates three independent evidence streams. The \emph{code},
via the inferrer of \cref{ch:inferrer} and \cref{ch:compositional}, answers what
the system does; the \emph{issue tracker} answers why, recording the defects and
edge cases the behaviour has provoked; and the \emph{documentation} answers
declared intent. The reconciliation is deliberately not a merge that smooths away
conflict. Its primary product is \emph{disagreement detection}: the points at
which the streams make incompatible claims are each a hypothesis that something is
wrong---the code, the documentation, or the recorded understanding. A clause the
code cannot satisfy but that a documentation- or ticket-derived stream asserts is
direct evidence of a code-versus-intent gap. This reframing forces a reordering of
the loop. In greenfield use it is prospective; the first brownfield pass is
\emph{retroactive}, resolving the recovered discrepancies and establishing a
ratified baseline before the pipeline is permitted to author any change. Because a
whole legacy system cannot be formally specified, recovery adopts a \emph{frontier}
strategy: it specifies only the module under change, the public API surface, and
the modules named by active tickets, treating the remainder as governed by weak
inherited-default contracts. Verification during this phase must never be read as a
proof of correctness; when the contract is dominated by code-inferred clauses,
verifying the code against it is nearly circular, establishing only that the
implementation is self-consistent. Correctness relative to intent comes from the
intent streams and, ultimately, from human ratification---never from the code
stream alone.

\subsection{Compositional change propagation as a lifecycle capability}\label{sec:int-propagation}

The final capability promotes the compositional refinement of
\cref{ch:compositional} from an inference-time accuracy gain to a lifecycle
operation. Because contracts compose along the call graph, a change to a callee's
contract mechanically \emph{recomputes} the proof obligations of every caller, and
the callers whose obligations can no longer be discharged are exactly the callers
the change breaks---identified by re-running the verifier, not the program. The
analysis is directional in a way that mirrors behavioural subtyping: weakening a
precondition or strengthening a postcondition is safe and flags no caller, whereas
strengthening a precondition or weakening a postcondition is potentially breaking
and yields a ranked, justified, per-call-site work list. The same machinery
answers two questions the prevailing tooling treats as unrelated---``does this
internal refactor break any call site?'' and ``does this dependency upgrade break
any call site?''---because, at the level of contracts, an internal refactor and a
third-party version bump are the same operation. This is the concrete link between
the two strands of the thesis core: formal-methods adoption and library
compatibility are served by one mechanism. It addresses a real gap, since semantic
versioning is an unreliable compatibility signal~\citep{preston2013semver,
raemaekers2017semver} and signature-level tools are blind to behavioural breaks
that leave the signature untouched~\citep{mostafa2017behavioral}. The carriage
mechanism of \cref{sec:int-carriage} is what makes the library side feasible
without source, and the cache-backed CI of \cref{sec:int-ci} is what makes
re-running the analysis on every upgrade economical. The soundness and the limits
are those of the underlying instrument: a caller reported compatible is compatible
over the whole specified domain, but a caller the solver cannot decide is reported
as undetermined rather than silently passed, and because an inferred callee
contract skews weak, a genuine break may be missed where the inferred contract
failed to capture the guarantee the caller relied upon. The analysis is therefore
sound in flagging breaks but not complete, and trustworthy only within the
specified frontier.

\section{Designed versus measured}\label{sec:int-disclaimer}

A scruple must be stated plainly and held throughout. This pipeline is a
\emph{reference architecture synthesised from validated components}, not an
empirically evaluated deployment. Its trusted core already exists and has been
independently validated in the earlier chapters: the JML-Inferrer and its
compositional pass that bootstrap draft contracts (\cref{ch:inferrer},
\cref{ch:compositional}), the OpenJML fork with Z3 as the oracle that never grades
itself, the artefact-embedding machinery of \cref{sec:int-carriage}, and the
Dockerised CI pipeline of \cref{sec:int-ci}. The only new software is untrusted
orchestration---sequencing the phases and managing the counterexample loop---plus
the model calls that draft and synthesise; none renders a verification verdict.
The architecture has been exercised only in part: the synthesis--verification loop
against fixed, known-dischargeable contracts on Apache Commons Lang, while contract
generation from intent and conversational ratification are specified as a staged
build plan with full validation deferred. The claim is accordingly one of
\emph{feasibility and realism}---that the spine can be assembled from parts that
already work and made to surface real code-versus-intent discrepancies on a genuine
three-stream corpus---not one of a measured effect at scale. The distinction
between a \emph{verified} obligation, which OpenJML has discharged, and a merely
\emph{well-formed} one, which is syntactically valid and scope-correct but
unproved, is never allowed to blur here: the compositional pass of
\cref{ch:compositional} contributes well-formed preconditions, and only those the
verifier discharges are verified. \Cref{ch:discussion} takes up the threats to this
position in full, and \cref{ch:conclusion} restates the boundary between what has
been built and what remains to be measured.
